\section{Related Works}
In this section, we provide an introduction to existing qLDPC decoders, most of which are developed based on the computational paradigms of belief propagation (BP) \cite{bp} or BP with ordered statistics decoding (BP+OSD) \cite{bposd}.

\subsection{BP-based qLDPC Decoders}
BP is a message-passing algorithm widely used in decoding qLDPC codes due to its linear-time complexity and suitability for sparse graph structures \cite{muller2025improved,miao2024joint}. However, standard BP suffers from low decoding accuracy due to its inability to fully account for quantum degeneracy and the presence of short cycles in the Tanner graph \cite{kuo2022exploiting}. To tackle this problem, numerous BP-based improved decoders have been proposed, such as enhanced feedback iterative decoding \cite{wang2012enhanced}, neural BP \cite{liu2019neural}, GF(2)-based augmented decoder \cite{rigby2019modified}, enhanced BP with trapping set dynamics \cite{chytas2024enhanced}. Recently, BP guided decimation (BP+GD) \cite{bpgd} has been proposed to encourage convergence by decimating the most reliable variable node in each iteration. Although these variants achieve higher accuracy than standard BP, their decoding performance still falls significantly short of that of maximum likelihood decoding (MLD) \cite{mld}, which remains the gold standard in terms of accuracy.

\subsection{Post-processing-based qLDPC Decoders}
In addition to BP-based improvements, some approaches enhance decoding accuracy through post-processing techniques, among which OSD is the most representative \cite{ninkovic2024decoding,bhattacharyya2025decoding}. However, OSD relies on computationally intensive operations, such as matrix inversion, which results in high decoding latency \cite{valls2021syndrome}. To solve this problem, BP+GDG \cite{bpgdg} proposes a sliding window decoder based on BP with guided decimation in the presence of circuit-level noise. Ambiguity clustering decoder \cite{wolanski2024ambiguity} divides the syndrome data into clusters that can be decoded independently, and decoders 144-qubit Gross code in 135$\mu s$ per round of syndrome extraction. Recently, BP+LSD \cite{bplsd} introduces a reliability-guided inversion decoder, which employs a parallel matrix factorization strategy to solve local decoding regions on the decoding graph. Although these methods improve scalability to some extent, they still involve matrix inversion operations, making it difficult to extend them to large-scale decoding problems.